%%%%%%%%%%%%%%%%%%%%%%%%%%%%%%%%%%%%%%%%%%%%%%%%%%%%%%%%%%%%%%%%%%%%%%
% Lớp 12 | Tọa độ của véc-tơ trong không gian | Tọa độ của véc-tơ và các công thức
% Dạng 2H2?2-5: Công thức tọa độ của tích có hướng và ứng dụng
% Số câu: 27  (sinh tự động bởi tach_cau.py; từ nay thầy sửa trực tiếp file này)
%%%%%%%%%%%%%%%%%%%%%%%%%%%%%%%%%%%%%%%%%%%%%%%%%%%%%%%%%%%%%%%%%%%%%%

%%%%%%%%%%%%%%%% TRẮC NGHIỆM 4 PHƯƠNG ÁN %%%%%%%%%%%%%%%%

% nguồn: Lop_12_Hinh_hoc_Chuong_2
\begin{ex}%[2H2N2-5]%[ma:f81c3ed2]
	 Trong không gian $Oxyz$, cho $\overrightarrow{a}=\left(x_0;y_0;z_0\right)$, $\overrightarrow{b}=\left(x_1;y_1;z_1\right)$. Tọa độ $\left[\overrightarrow{a},\overrightarrow{b}\right]$ là
	 \choice
		{$\left(y_0z_1-y_1z_0;x_0z_1-x_1z_0;x_0y_1-x_1y_0\right)$} 
		{\True $\left(y_0z_1-y_1z_0;-x_0z_1+x_1z_0;x_0y_1-x_1y_0\right)$} 
		 {$\left(y_0z_1+y_1z_0;x_0z_1+x_1z_0;x_0y_1+x_1y_0\right)$}
		{$\left(y_0z_1-y_1z_0;-x_0z_1-x_1z_0;x_0y_1-x_1y_0\right)$} 
	 \loigiai{
		Tọa độ  $\left[\overrightarrow{a},\overrightarrow{b}\right]=\left(y_0z_1-y_1z_0;-x_0z_1+x_1z_0;x_0y_1-x_1y_0\right)$.
	} 
\end{ex}

% nguồn: Lop_12_Hinh_hoc_Chuong_2
\begin{ex}%[2H2N2-5]%[ma:458b5c45]
	Trong không gian với hệ tọa độ $Oxyz$, cho hai vectơ $\overrightarrow{a} = (a_1;a_2;a_3)$,\break $\overrightarrow{b} = (b_1;b_2;b_3)$ khác $\overrightarrow{0}$. Tích  có hướng của $\overrightarrow{a}$ và $\overrightarrow{b}$ là $\overrightarrow{c}$. Khẳng định nào sau đây đúng?
	\choice 
		{$\overrightarrow{c} = (a_1b_2  - a_2b_1, a_2b_3 - a_3b_2, a_3b_1 - a_1b_3)$}
		{$\overrightarrow{c} = (a_1b_2 - a_2b_1, a_3b_2 - a_2b_3,  a_3b_1 - a_1b_3)$}
		{$\overrightarrow{c} = (a_1b_2 - a_2b_1, a_3b_2 - a_2b_3, a_1b_3  - a_3b_1)$}
		{\True $\overrightarrow{c} = (a_2b_3 - a_3b_2,  a_3b_1 - a_1b_3, a_1b_2 - a_2b_1)$}
	 \loigiai{
		Ta có $\overrightarrow{c} =  [\overrightarrow{a}, \overrightarrow{b}] = (a_2b_3 - a_3b_2, a_3b_1 - a_1b_3, a_1b_2 - a_2b_1)$.
	} 
\end{ex}

% nguồn: Lop_12_Hinh_hoc_Chuong_2
\begin{ex}%[2H2N2-5]%[ma:c34ed409]
	Trong không gian $Oxyz$, cho hai vectơ $\vec{u}=(3; 1;  -2)$ và $\vec{v}=(-2; 1; 5)$. Toạ độ của vectơ $\left[\vec{u}, \vec{v}\right]$~là
	 \choice
		{$(5;  7; -11)$}
		{$(-7; 11;  -5)$}
		{\True $(7; -11; 5)$} 
		{$(-5; -7;  11)$}
	\loigiai{ 
		Tích có hướng  của hai vectơ $\vec{u}$ và $\vec{v}$~là
		\[\left[\vec{u}, \vec{v}\right]= \left(1\cdot5-(-2)\cdot 1;(-2)\cdot(-2)-3\cdot 5; 3\cdot 1-1\cdot(-2)  \right)=(7; -11; 5).\]
	 }
\end{ex}

% nguồn: Lop_12_Hinh_hoc_Chuong_2
\begin{ex}%[2H2N2-5]%[ma:df7c2e87]
	Trong không gian với hệ tọa độ $Oxyz$, cho hai  vectơ $\overrightarrow{a} = (2;1;-2)$ và vectơ $\overrightarrow{b} = (1;0;2)$. Tìm tọa độ vectơ $\overrightarrow{c}$ là tích có hướng của $\overrightarrow{a}$ và $\overrightarrow{b}$.
	\choice 
		 {$\overrightarrow{c} = (2;6;-1)$}
		{$\overrightarrow{c} =  (4;6;-1)$}
		{$\overrightarrow{c}  = (4;-6;-1)$}
		{\True $\overrightarrow{c}  = (2;-6;-1)$}
	\loigiai{ 
		Công thức tính tích có hướng trong hệ trục  tọa độ $Oxyz$ ta được
		$$\overrightarrow{c} =  [\overrightarrow{a}, \overrightarrow{b}]= (2; -6; -1).$$
	} 
\end{ex}

% nguồn: Lop_12_Hinh_hoc_Chuong_2
\begin{ex}%[2H2N2-5]%[ma:f8a6ec56]
	 Trong không gian $Oxyz$, cho $\overrightarrow{u}=(1;2;3)$, $\overrightarrow{v}=(0;-1;1)$. Tìm tọa độ của vectơ tích có hướng của hai vectơ $\overrightarrow{u}$ và $\overrightarrow{v}$.
	\choice 
		 {$(5;1;-1)$}
		{\True  $(5;-1;-1)$}
		{$(-1;-1;-1)$} 
		 {$(-1;-1;5)$}
	 \loigiai{
		Tích có hướng của hai vectơ  $\overrightarrow{u}$ và $\overrightarrow{v}$ là $\left[\overrightarrow{u}, \overrightarrow{v}\right] =(5;-1;-1)$.
	} 
\end{ex}

% nguồn: Lop_12_Hinh_hoc_Chuong_2
\begin{ex}%[2H2H2-5]%[ma:0eb574ec]
	Trong không gian $Oxyz$, cho hai vectơ $\overrightarrow{m}=(4;3;1)$ và $\overrightarrow{n}=(0;0;1)$. Gọi $\overrightarrow{p}$ là vectơ cùng hướng với $\left[\overrightarrow{m}, \overrightarrow{n}\right]$ và $\left| \overrightarrow{p}\right|=15$. Tọa độ của vectơ  $\overrightarrow{p}$ là
	 \choice
		 {$(0;9;-12)$}
		{$(-9;12;0)$} 
		 {$(0;-9;12)$}
		 {\True $(9;-12;0)$}
	 \loigiai{
		Ta có $\left[\overrightarrow{m},  \overrightarrow{n}\right]=(3;-4;0)$.\\
		Vì $\overrightarrow{p}$ là vectơ  cùng hướng với $\left[\overrightarrow{m}, \overrightarrow{n}\right]$ nên $\overrightarrow{p}=k(3;-4;0)$.\\
		Hơn nữa $\left|\overrightarrow{p}\right|=15 \Leftrightarrow k \sqrt{3^2+4^2}=15 \Leftrightarrow k=3$.  \\
		Vậy $\overrightarrow{p}=(9;-12;0)$. 
	 }
\end{ex}

% nguồn: Lop_12_Hinh_hoc_Chuong_2
\begin{ex}%[2H2H2-5]%[ma:c23da12f]
	Trong không gian $Oxyz$, cho tam giác $ABC$ với $A(1;2;3)$, $B(0;1;4)$ và  $C(2;3;-2)$. Tính diện tích $S$ của tam giác $ABC$.
	 \choice
		{\True $S=2  \sqrt{2}$}
		 {$S=6 \sqrt{2}$}
		{$S=4 \sqrt{2}$} 
		 {$S=3 \sqrt{2}$}
	\loigiai{ 
		 Ta có $\overrightarrow{AB}=(-1;-1;1)$, $\overrightarrow{AC}=(1;1;-5)$. \\
		Vậy diện tích $S$ của tam giác $ABC$ là $S=\dfrac{1}{2}\left|\left[\overrightarrow{AB},\overrightarrow{AC}\right]\right|=2  \sqrt{2}$.
	} 
\end{ex}

% nguồn: Lop_12_Hinh_hoc_Chuong_2
\begin{ex}%[2H2H2-5]%[ma:8b026ebf]
	Trong không gian $Oxyz$, cho $\overrightarrow{a}=(-3;-1;1)$, $\overrightarrow{b}=(4;1;2)$, $\overrightarrow{c}=(1;0;m+2)$. Tìm $m$ để ba vectơ  $\overrightarrow{a}$, $\overrightarrow{b}$, $\overrightarrow{c}$ đồng phẳng.
	\choice 
		{$m=-5$} 
		{$m=5$} 
		 {$m=-1$}
		 {\True $m=1$}
	\loigiai{ 
		Ta có $\left[\overrightarrow{a}, \overrightarrow{b}\right] =\left(\left|\begin{array}{cc}-1 & 1 \\ 1 & 2
	\end{array}\right|;\left|\begin{array}{cc}1 & -3 \\ 2 & 4
	\end{array}\right|;\left|\begin{array}{cc}-3 & -1 \\ 4 & 1
	\end{array}\right|\right)=(-3;10;1)$.\\
	Mà $\left[\overrightarrow{a}, \overrightarrow{b}\right]\cdot \overrightarrow{c}=(-3) \cdot  1+1 \cdot(m+2)=m-1$.\\
	Ba vectơ  $\overrightarrow{a}$, $\overrightarrow{b}$, $\overrightarrow{c}$ đồng phẳng $\Leftrightarrow\left[\overrightarrow{a}, \overrightarrow{b}\right]\cdot \overrightarrow{c}=0 \Leftrightarrow m-1=0 \Leftrightarrow m=1$.
} 
\end{ex}

% nguồn: Lop_12_Hinh_hoc_Chuong_2
\begin{ex}%[2H2H2-5]%[ma:ae1366ab]
	Trong không gian $Oxyz$, cho $A(1;-2;0)$, $B(1;0;-1)$, $C(0;-1;2)$ và $D(0;3; m)$. Giá trị của $m$  thuộc khoảng nào sau đây để bốn điểm trên đồng phẳng?
	 \choice
		 {$(-2;-1)$}
		 {\True $(-1;1)$}
		{$(1;2)$} 
		 {$(5;7)$}
	\loigiai{ 
		Ta  có $\overrightarrow{AB}=(0;2;-1)$, $\overrightarrow{AC}=(-1;1;2)$, $\overrightarrow{AD}=(-1;5;m)$, $\left[\overrightarrow{AB},\overrightarrow{AC}\right]=(5;1;2)$.\\
		Lúc đó $A$, $B$, $C$, $D$ đồng phẳng khi và chỉ khi $\left[\overrightarrow{AB},\overrightarrow{AC}\right] \cdot  \overrightarrow{AD}=0 \Leftrightarrow 2m=0 \Leftrightarrow m=0$.
	 }
\end{ex}

% nguồn: Lop_12_Hinh_hoc_Chuong_2
\begin{ex}%[2H2H2-5]%[ma:0147c8f5]
	Trong không gian $Oxyz$, cho hai điểm  $A(1;2;3)$, $B(3;4;5)$. Diện tích tam giác $OAB$ bằng
	 \choice
		{$\sqrt{87}$} 
		{$\dfrac{\sqrt{21}}{2}$} 
		 {\True $\sqrt{6}$}
		{$2  \sqrt{6}$}
	 \loigiai{
		Ta có  $\overrightarrow{OA}=(1;2;3)$, $\overrightarrow{OB}=(3;4;5)$, $\left[\overrightarrow{OA},\overrightarrow{OB}\right]=(-2;4;-2)$.\\
		Lúc  đó
		 $S_{\triangle OAB}=\dfrac{1}{2}\left|\left[\overrightarrow{OA},\overrightarrow{OB}\right]\right|=\dfrac{1}{2} \sqrt{(-2)^2+4^2+(-2)^2}=\sqrt{6}$.\\
		Vậy diện tích tam giác $OAB$ là $\sqrt{6}$. 
	} 
\end{ex}

% nguồn: Lop_12_Hinh_hoc_Chuong_2
\begin{ex}%[2H2H2-5]%[ma:d91d1b18]
	Trong không  gian $Oxyz$, cho điểm $A(-4;6;2)$. Gọi $M$, $N$, $P$ lần lượt là hình chiếu của $A$ trên các trục $Ox$, $Oy$ và $Oz$. Tính diện tích $S$ của tam giác $MNP$.
	 \choice
		{$S=28$} 
		{$S=\dfrac{49}{2}$} 
		 {$S=7$}
		{\True  $S=14$}
	\loigiai{ 
		Theo đề bài ta  có $M(-4;0;0)$, $N(0;6;0)$ và $N(0;0;2)$, $\left[\overrightarrow{MN}, \overrightarrow{MP}\right] =(12;-8;-24)$.\\
		Diện tích của tam giác $MNP$ là $$S=\dfrac{1}{2}\left|\left[\overrightarrow{MP}, \overrightarrow{MN}\right]\right|=\dfrac{1}{2}  \sqrt{12^2+(-8)^2+(-24)^2}=14.$$
		 Vậy $S=14$.
	 }
\end{ex}

% nguồn: Lop_12_Hinh_hoc_Chuong_2
\begin{ex}%[2H2V2-5]%[ma:7939080e]
	Cho hình lập phương $ABCD.A'B'C'D'$ có độ dài cạnh bằng $1$. Gọi $M$, $N$, $P$, $Q$, lần lượt là trung điểm của các cạnh $AB$, $BC$, $C'D'$ và $DD'$. Thể tích khối tứ diện $MNPQ$ là 
	\choice 
		 {\True $\dfrac{1}{12}$}
		{$\dfrac{1}{8}$} 
		{$\dfrac{3}{8}$} 
		{$\dfrac{1}{24}$} 
	 \loigiai{
		Chọn hệ trục tọa độ $Oxyz$ sao cho:  $D'\equiv O$; $Ox \equiv D'A'$; $Oy\equiv D'C'$; $Oz\equiv D'D$.\\
		Khi đó $A(1;0;1)$, $B(1;1;1)$, $C(0;1;1)$, $D(0;0;1)$, $A'(1;0;0)$, $B'(1;1;0)$,  $C'(0;1;0)$.
		\begin{center}
			\begin{tikzpicture}[line join = round,line cap=round,>=stealth,font=\footnotesize,scale=1]
					\path 
						(1,1) coordinate (A') 
						(1,4.5) coordinate  (A)
						(5.25,1) coordinate  (B')
						(5.25,4.5) coordinate (B) 
						(3.95,2.45) coordinate (C') 
						(3.95,5.9) coordinate (C) 
						 (-0.25,2.45) coordinate (D')
						(-0.25,5.9) coordinate  (D)
						($(B)! 1/2 !(A)$)  coordinate (M)
						($(B)!  1/2 !(C)$) coordinate (N)
						($(D)!  1/2 !(D')$) coordinate (Q)
						 ($(D')! 1/2 !(C')$) coordinate (P);
					 \foreach \p/\r in {A/90,B/-130,C/0,D/180,A'/180,B'/0,C'/45,D'/180,M/-70,N/0,P/-90,Q/180}
					\fill (\p) circle (1.2pt) node[shift={(\r:3mm)}]{$\p$}; 
					\draw (D)--(D') 
						 (A)--(A')
						(D')--(A')--(B') 
						 (C)--(D)--(A)--(B)--(B')
						(B)--(C) 
						(M)--(N); 
					\draw [dashed]  (C)--(C')--(B')
						 (A')--(D')--(C')
						 (N)--(P)
						 (Q)--(M)
						(Q)--(N) 
						 (Q)--(P)--(N)
						(P)--(M); 
					\draw[->] (D)--++(0,1) node[right]{$z$}; \draw[->] (C')--++(1,0)  node[above]{$y$}; \draw[->] (A')--++(1,-1) node[left]{$x$};
				\end{tikzpicture}
		\end{center}
		Ta có  $M\left(1;\dfrac{1}{2};1\right)$, $N\left(\dfrac{1}{2};1;1\right)$, $P\left(0;\dfrac{1}{2};0\right)$, $Q\left(0;0;\dfrac{1}{2}\right)$.\\
		Ta có $\overrightarrow{MN}=\left(-\dfrac{1}{2}; \dfrac{1}{2};  0\right)$, $\overrightarrow{MP}=(-1; 0; -1)$, $\overrightarrow{MQ}=\left(-1; -\dfrac{1}{2}; -\dfrac{1}{2}\right)$.\\
		$\left[\overrightarrow{MN},  \overrightarrow{MP}\right]= \left(-\dfrac{1}{2}, -\dfrac{1}{2}, \dfrac{1}{2}\right)$.\\
		$\left|\left[\overrightarrow{MN}, \overrightarrow{MP}\right]\cdot\overrightarrow{MQ}\right| = \left|\dfrac{1}{2} + \dfrac{1}{4} - \dfrac{1}{4}\right| = \dfrac{1}{2} \Rightarrow V_{MNPQ} = \dfrac{1}{6}\left|\left[\overrightarrow{MN},  \overrightarrow{MP}\right]\cdot\overrightarrow{MQ}\right| = \dfrac{1}{12}$.
	 }
\end{ex}

% nguồn: Lop_12_Hinh_hoc_Chuong_2
\begin{ex}%[2H2V2-5]%[ma:8fbbc63a]
	Trong không gian $Oxyz$, cho hình chóp $S.ABCD$ có đáy $ABCD$ là hình thoi, $AC$ cắt $BD$ tại gốc tọa độ $O$. Biết $A(2;0;0)$, $B(0;1;0)$, $S(0;0;2\sqrt{2})$. Gọi $M$ là trung điểm của cạnh  $SC$. Mặt phẳng $(ABM)$ cắt đường thẳng $SD$ tại điểm $N$. Tính thể tích khối chóp $S.ABMN$.
	 \choice
		{\True  $V = \sqrt{2}$}
		{$V = 2\sqrt{3}$} 
		 {$V = 3\sqrt{2}$}
		{$V = \dfrac{\sqrt{3}}{4}$} 
	 \loigiai{
		\begin{center}
			\begin{tikzpicture}[line join = round,line cap=round,>=stealth,font=\footnotesize,scale=1]
					\path 
						(0,0) coordinate  (A)
						(2,1)  coordinate (D)
						 (5.5,1) coordinate (C)
						(3.5,0) coordinate (B) 
						(2.75,5)  coordinate (S)
						(intersection of B--D and  A--C) coordinate (O)
						 ($(S)! 1/2 !(C)$) coordinate (M)
						 ($(S)! 1/2 !(D)$) coordinate (N);
					 \draw (A)--(B)--(C)
						 (S)--(A)
						(S)--(B) 
						 (S)--(C);
					 \draw[dashed] (A)--(D)--(C)
						(O)--(S)--(D) 
						(D)--(B) 
						 (A)--(C)
						(A)--(M)--(B) 
						 (A)--(N)--(M);
					\draw[->] (S)--++(0,1) node[right]{$z$}; \draw[->] (A)--++(-1,-0.2) node[below]{$x$}; \draw[->]  (B)--++(1,-0.7) node[right]{$y$};
					\foreach  \p/\r in {A/-90,D/150,C/0,B/-90,S/130,O/-90,M/45,N/145}
					\fill (\p) circle (1.2pt)  node[shift={(\r:3mm)}]{$\p$};
				\end{tikzpicture}
		\end{center}
		Ta có $A(2;0;0)$,  $B(0;1;0)$, $S(0;0;2\sqrt{2})$, suy ra tọa độ các điểm $D(0;-1;0)$, $C(-2;0;0)$, $M\left( -1;0;\sqrt{2}\right) $.\\
		Ta có $MN\parallel AB\parallel CD$ $\Rightarrow N$ là trung điểm $SD$ $\Rightarrow N\left(0;  -\dfrac{1}{2};\sqrt{2}\right)$.\\
		$\overrightarrow{SA}= \left(2;0;-2\sqrt{2}\right)$, $\overrightarrow{SM}=\left(-1;0;-\sqrt{2}\right)$,  $\overrightarrow{SB}=\left(0;1;-2\sqrt{2}\right)$, $\overrightarrow{SN}=\left(0;-\dfrac{1}{2};-\sqrt{2}\right)$.\\
		$\Rightarrow \left[\overrightarrow{SA},\overrightarrow{SM}\right]=\left(0;4\sqrt{2};0\right)$.\\ 
		Vậy  thể tích chóp $S.ABMN$ là\\
		$V = V_{S.ABM} + V_{S.AMN} = \dfrac{1}{6}\left|\left[\overrightarrow{SA}, \overrightarrow{SM}\right]\cdot\overrightarrow{SB}\right| + \dfrac{1}{6}\left|\left[\overrightarrow{SA}, \overrightarrow{SM}\right]\cdot\overrightarrow{SN}\right| = \dfrac{2\sqrt{2}}{3} + \dfrac{\sqrt{2}}{3}  = \sqrt{2}$.
	} 
\end{ex}

% nguồn: Lop_12_Hinh_hoc_Chuong_2
\begin{ex}%[2H2V2-5]%[ma:7b16f2de]
	Trong không gian với hệ trục $Oxyz$, cho hình hộp $ABCD.A'B'C'D'$. Biết rằng $\overrightarrow{A B}=(1;3;4)$, $\overrightarrow{A D}=(-2;3;5)$ và $\overrightarrow{A C'}=(1;1;1)$. Tính thể tích hình hộp  $ABCD.A'B'C'D'$.
	 \choice
		{$V_{ABCD.A'B'C'D'}=6$} 
		 {$V_{ABCD.A'B'C'D'}=12$}
		{\True $V_{ABCD.A'B'C'D'}=1$} 
		{$V_{ABCD \cdot A'B'C'D'}=3$} 
	\loigiai{Ta có $V_{ABCD.A'B'C'D'}=\left|\left[\overrightarrow{A B}, \overrightarrow{A D}\right] \cdot \overrightarrow{A A'}\right|$, trong đó  $\left[\overrightarrow{A B}, \overrightarrow{A D}\right]=(3;-13;9)$.\\
		Lại có $\overrightarrow{A C'}=\overrightarrow{A B}+\overrightarrow{A D}+\overrightarrow{A  A'} \Rightarrow \overrightarrow{A A'}=\overrightarrow{A C'}-\left(\overrightarrow{A B}+\overrightarrow{A D}\right)$.
		Suy  ra $\overrightarrow{A A'}=(2;-5;-8)$.\\
		Vậy  $V_{ABCD.A'B'C'D'}=\left|\left[\overrightarrow{A B}, \overrightarrow{A D}\right] \cdot \overrightarrow{A A'}\right|=|6+65-72|=1$.
	} 
\end{ex}

% nguồn: Lop_12_Hinh_hoc_Chuong_2
\begin{ex}%[2H2V2-5]%[ma:c3799b42]
	Cho hình chóp $S.ABC$ có đáy $ABC$ là tam giác vuông cân tại $A$ và $AB = a$, cạnh bên $SA$ vuông góc  với mặt phẳng đáy và $SA = a$. Gọi $D$, $E$, $F$ lần lượt là điểm đối xứng của $A$ qua $C$, của $S$ qua $B$ và của $A$ qua mặt phẳng $(SBC)$. Thể tích của khối tứ diện $ADEF$ bằng
	\choice 
		{\True  $\dfrac{2a^3}{3}$}
		{$\dfrac{2\sqrt{3}a^3}{3}$} 
		{$\dfrac{a^3}{3}$} 
		{$\dfrac{\sqrt{3}a^3}{3}$} 
	 \loigiai{
		\begin{center}
			\begin{tikzpicture}[line join = round,line cap=round,>=stealth,font=\footnotesize,scale=1]
					 \path
						(0,0.7) coordinate  (A)
						(2.1,0) coordinate  (B)
						 (3.55,0.7) coordinate (C)
						 (0,4.65) coordinate (S)
						(1.9,1.76) coordinate (H) 
						(7.1,0.7)  coordinate (D)
						(4.22,-4.65)  coordinate (E)
						(3.78,2.85) coordinate  (F)
						(2.35,2.03) coordinate (N) 
						($(B)! 1/2 !(C)$)  coordinate (M);
					\draw (S)--(A)--(B)--cycle (H)--(F)--(D)--(E)--(A) 
						 (F)--(E)
						 (H)--(S)--(N)
						(S)--(E); 
					\draw [dashed]  (A)--(C)--(B)
						 (C)--(D)
						(A)--(H)--(M) 
						(N)--(C)--(H); 
					\draw[->] (S)--++(0,1) node[left]{$z$}; \draw[->] (D)--++(1,0) node[below]{$y$};  \draw[->] (B)--++(0.7,-0.25) node[below]{$x$};
					\draw pic[draw,angle radius=3mm] {right angle = A--H--M}; 
					\draw pic[draw,angle  radius=3mm] {right angle = S--A--B};
					\draw pic[draw,angle radius=3mm] {right  angle = S--A--C};
					 \draw pic[draw,angle radius=3mm] {right angle = B--A--C};
					 \foreach \p/\r in {A/-90,B/-90,C/-50,S/180,H/90,D/-50,E/-90,F/90}
					\fill (\p) circle (1.2pt)  node[shift={(\r:3mm)}]{$\p$};
				\end{tikzpicture}
		\end{center}
		Chọn hệ trục tọa độ $Axyz$ như hình vẽ  với độ dài vectơ đơn vị bằng $a$, chọn $a=1$ ta được tọa độ các điểm như sau
		$A(0;0;0)$, $B(1;0;0)$,  $C(0;1;0)$, $S(0;0;1) \Rightarrow D(0;2;0)$, $E(2;0;-1)$.\\
		Gọi $H$ là trọng tâm tam giác $SBC \Rightarrow H\left(\dfrac{1}{3}; \dfrac{1}{3}; \dfrac{1}{3}\right)$. Vì tam giác $SBC$ là tam giác đều nên $H$ cũng là trực tâm của tam giác $SBC \Rightarrow AH \perp (SBC) \Rightarrow F$ là điểm đối  xứng của $A$ qua $H \Rightarrow F\left(\dfrac{2}{3}; \dfrac{2}{3}; \dfrac{2}{3}\right)$.\\
		Ta tính được $\overrightarrow{AD}=(0;2;0)$, $\overrightarrow{AE}=(2;0;-1)$, $\overrightarrow{AF}=\left(\dfrac{2}{3};\dfrac{2}{3};\dfrac{2}{3}\right)  \Rightarrow \left[\overrightarrow{AD}, \overrightarrow{AE}\right] = (-2;0;-4)$.\\
		Khi đó $V_{A.DEF} = \dfrac{1}{6}\left|\left[\overrightarrow{AD}, \overrightarrow{AE}\right]\cdot\overrightarrow{AF}\right| =  \dfrac{1}{6}\left|-2\cdot\dfrac{2}{3} + 0\cdot\dfrac{2}{3} - 4\cdot\dfrac{2}{3}\right| = \dfrac{2}{3}$.\\
		Vì độ dài vectơ đơn vị trên trục  bằng $a$ nên ta được thể tích cần tìm là $\dfrac{2a^3}{3}$.
	 }
\end{ex}

% nguồn: Lop_12_Hinh_hoc_Chuong_2
\begin{ex}%[2H2V2-5]%[ma:e61371d3]
	Trong không gian với hệ tọa độ $Oxyz$, cho bốn điểm $A(1;1;4)$, $ B(5;-1;3)$, $C(3;1;5)$ và $D(2;2;m)$ (với $m$ là tham số). Xác định $m$ để bốn điểm $A$, $B$, $C$, $D$ tạo thành bốn đỉnh của một tứ  diện.
	\choice 
		{\True  $m \neq 6$}
		{$m  \neq 4$}
		{$m \in  \mathbb{R}$}
		 {$m<0$}
	\loigiai{Ta  có
		 \allowdisplaybreaks
		\begin{eqnarray*}
			\overrightarrow{A B}&=&(4;-2;-1), \overrightarrow{A C}=(2;0;1), \overrightarrow{A D}=(1;1;m-4).\\ 
			\left[\overrightarrow{A B}, \overrightarrow{A C}\right]&=&(-2;-6;4).\\\left[\overrightarrow{A B}, \overrightarrow{A C}\right] \cdot \overrightarrow{A  D}&=&-2-6+4(m-4)=4 m-24.
		\end{eqnarray*}
		Bốn điểm $A$, $B$, $C$, $D$ tạo thành bốn đỉnh của một tứ diện  khi và chỉ khi
		$$\left[\overrightarrow{A B}, \overrightarrow{A  C}\right] \cdot \overrightarrow{A D}=4 m-24 \neq 0 \Leftrightarrow m \neq 6.$$
	} 
\end{ex}

% nguồn: Lop_12_Hinh_hoc_Chuong_2
\begin{ex}%[2H2V2-5]%[ma:3a3ca03f]
	Trong không gian với hệ tọa độ $Oxyz$, cho hình hộp chữ nhật $ABCD.A'B'C'D'$ có $A$ trùng với gốc tọa độ $O$. Biết rằng $B(m;0;0)$, $D(0;m;0)$, $A'(0;0;n)$ với $m, n$ là các số dương và $m+n=4$. Gọi  $M$ là trung điểm của cạnh $CC'$. Thể tích lớn nhất của khối tứ diện $BDA'M$ bằng
	\choice 
		 {\True $\dfrac{64}{27}$}
		{$\dfrac{11\sqrt{2}}{6}$} 
		{$\dfrac{9}{4}$} 
		{$\dfrac{75}{32}$} 
	\loigiai{ 
		\begin{center}
			\begin{tikzpicture}[line join = round,line cap=round,>=stealth,font=\footnotesize,scale=1]
					\path 
						(0,0)  coordinate (B)
						(1.25,1.5) coordinate (A) 
						(4,0) coordinate  (C)
						(5.25,1.5) coordinate  (D)
						(0,4) coordinate  (B')
						(1.25,5.5) coordinate  (A')
						(4,4) coordinate (C') 
						(5.25,5.5)  coordinate (D')
						($(C)! 1/2 !(C')$)  coordinate (M);
					 \draw (D)--(C)--(B)
						(A')--(B')--(C')--(D')--cycle (D)--(D') 
						 (B)--(B')
						 (C)--(C');
					\draw [dashed]  (D)--(A)--(B)
						(A)--(A') 
						 (A')--(D)--(B)--cycle (M)--(A')
						(M)--(D) 
						(M)--(B); 
					\draw[->] (B)--++(-0.8,-1)  node[left]{$x$};]
					\draw[->] (A')--++(0,1)  node[right]{$z$};
					\draw[->] (D)--++(1,0)  node[above]{$y$};
					\foreach  \p/\r in {D'/90,C'/-50,B'/90,A'/150,D/-90,C/-50,B/-50,A/180,M/180}
					\fill (\p) circle  (1.2pt) node[shift={(\r:3mm)}]{$\p$};
				\end{tikzpicture}
		\end{center}
		Ta có $A(0;0;0)$,  $B(m;0;0)$, $D(0;m;0)$, $A'(0;0;n)$.\\
		 Suy ra $C(m;m;0)$, $B'(m;0;n)$, $C'(m;m;n)$, $D'(0;m;n)$, $M\left(m;m;\dfrac{n}{2}\right)$.\\
		Và $\overrightarrow{BD} = (-m;m;0)$, $\overrightarrow{BA'} = (-m;0;n)$,  $\overrightarrow{BM} = \left(0;m;\dfrac{n}{2}\right)$.\\
		$\left[\overrightarrow{BD}, \overrightarrow{BA'}\right] = (mn; mn; m^2)$;  $\left[\overrightarrow{BD}, \overrightarrow{BA'}\right]\cdot\overrightarrow{BM} = \dfrac{3}{2}m^2n$.\\
		 Ta được
		$V_{BDA'M} = \dfrac{1}{6}\left|\left[\overrightarrow{BD},  \overrightarrow{BA'}\right]\cdot\overrightarrow{BM}\right| = \dfrac{1}{4}|m^2\cdot n| = \dfrac{1}{4}|m^2\cdot(4-m)|=\left|m^2-\dfrac{1}{4}m^3 \right|$.\\
		Xét $f(m)=m^2-\dfrac{1}{4}m^3$, ta  có $f'(m)=2m-\dfrac{3}{4}m^2$.\\
		Với $f'(m)=0\Leftrightarrow \hoac{&m=0 \text{ (không  thỏa số dương)}\\&m=\dfrac{8}{3} \text{ (thỏa số dương)}.}$\\
		 Ta có bảng biến thiên xét $m\in(0;4)$
		\begin{center}
			\begin{tikzpicture}
					 \tkzTabInit[nocadre=false,lgt=1.4,espcl=2.8,deltacl=0.5]
						{$m$/0.7,$f'(m)$/0.7,$f  (m)$/2}
						{$0$,$\frac{8}{3}$,$4$} 
					\tkzTabLine{,+,0,-,} 
					\tkzTabVar{-/$0$,+ /$\dfrac{64}{27}$,-/$0$} 
				\end{tikzpicture}
		\end{center}
		Do đó giá trị lớn nhất của $f(m)$ trên $(0;4)$ là $\dfrac{64}{27}$, dựa vào tính đơn điệu của $f(m)$, ta có  giá trị lớn nhất của hàm số $y=|f(m)|$ là $\dfrac{64}{27}$.\\
		Thể tích lớn nhất của  khối tứ diện $BDA'M$ bằng $\dfrac{64}{27}$.
	} 
\end{ex}

% nguồn: Lop_12_Hinh_hoc_Chuong_2
\begin{ex}%[2H2V2-5]%[ma:13f4acfb]
	Cho hình lập phương $ABCD.A'B'C'D'$ có độ dài cạnh bằng $1$. Gọi $M$, $N$,  $P$, $Q$ lần lượt là trung điểm của các cạnh $AB$, $BC$, $C'D'$ và $DD'$. Tính thể tích khối tứ diện $MNPQ$.
	\choice 
		{\True  $\dfrac{1}{12}$}
		 {$\dfrac{1}{8}$}
		 {$\dfrac{3}{8}$}
		{$\dfrac{1}{24}$} 
	 \loigiai{
		\begin{center}
			\begin{tikzpicture}[line join = round,line cap=round,>=stealth,font=\footnotesize,scale=1]
					 \path
						(0,0) coordinate (A) 
						(1.25,1.5) coordinate (D) 
						(4,0)  coordinate (B)
						(5.25,1.5) coordinate (C) 
						(0,4) coordinate  (A')
						(1.25,5.5) coordinate (D') 
						(4,4) coordinate  (B')
						(5.25,5.5)  coordinate (C')
						 ($(A)! 1/2 !(B)$) coordinate (M)
						($(B)! 1/2 !(C)$) coordinate (N) 
						($(C')! 1/2 !(D')$)  coordinate (P)
						($(D)! 1/2 !(D')$)  coordinate (Q);
					\draw  (A)--(B)--(C)
						 (A')--(B')--(C')--(D')--cycle (A)--(A')
						 (B)--(B')
						(C)--(C'); 
					 \draw [dashed] (A)--(D)--(C)
						(D)--(D') 
						 (M)--(N)--(P)--(Q)--cycle (Q)--(N)
						 (M)--(P);
					\draw[->] (A)--++(-0.8,-1) node[left]{$x$}; \draw[->] (A)--++(-0.8,-1) node[left]{$x$}; \draw[->] (D')--++(0,1) node[right]{$z$}; \draw[->]  (C)--++(1,0) node[above]{$y$};
					\foreach \p/\r in {A'/180,B'/-50,C'/90,D'/150,A/-90,B/-50,C/-50,D/180,M/-90,N/0,Q/180,P/90} 
					\fill  (\p) circle (1.2pt) node[shift={(\r:3mm)}]{$\p$};
				\end{tikzpicture}
		\end{center}
		Gắn hệ trục tọa độ như hình vẽ ta có $D(0;0;0)$, $A(1;0;0)$,  $B(1;1;0)$, $C(0;1;0)$, $A'(1;0;1)$, $B'(1;1;1)$, $C'(0;1;1)$, $D'(0;0;1)$.\\
		Vì $M$, $N$, $P$, $Q$ lần lượt là  trung điểm của các cạnh $AB, BC, C'D'$ và $DD'$ nên\\
		 $M\left(1;\dfrac{1}{2};0\right)$, $N\left(\dfrac{1}{2};1;0\right)$, $P\left(0;\dfrac{1}{2};1\right)$, $Q\left(0;0;\dfrac{1}{2}\right)$.\\
		Suy ra $\overrightarrow{MN} = \left(-\dfrac{1}{2}; \dfrac{1}{2}; 0\right)$, $\overrightarrow{MP}  = (-1; 0; 1)$, $\overrightarrow{MQ} = \left(-1; -\dfrac{1}{2}; \dfrac{1}{2}\right)$.\\
		Suy ra $\left[\overrightarrow{MN}, \overrightarrow{MP}\right] = \left(\dfrac{1}{2}; \dfrac{1}{2}; \dfrac{1}{2}\right) \Rightarrow \left[\overrightarrow{MN}, \overrightarrow{MP}\right]\cdot\overrightarrow{MQ} = -\dfrac{1}{2}$.\\ 
		Vậy $V_{MNPQ} = \dfrac{1}{6}\left|\left[\overrightarrow{MN}, \overrightarrow{MP}\right]\cdot\overrightarrow{MQ}\right| = \dfrac{1}{6}\left|-\dfrac{1}{2}\right| =  \dfrac{1}{12}$.
	 }
\end{ex}

% nguồn: Lop_12_Hinh_hoc_Chuong_2
\begin{ex}%[2H2V2-5]%[ma:062b0b9f]
	Trong không gian với hệ tọa độ $Oxyz$, cho hình thang $A B C D$ có hai đáy $A B$, $C  D$; có tọa độ ba đỉnh $A(1;2;1)$, $B(2;0;-1)$, $C(6;1;0)$. Biết hình thang có diện tích bằng $6 \sqrt{2}$. Giả sử đỉnh $D(a;b;c)$, tìm mệnh đề đúng?
	\choice 
		 {$a+b+c=6$}
		 {$a+b+c=5$}
		{\True $a+b+c=8$} 
		{$a+b+c=7$} 
	\loigiai{Ta có 
		$\overrightarrow{A  B}=(1;-2;-2)$, $\overrightarrow{D C}=(6-a;1-b;-c)$, $\overrightarrow{A C}=(5;-1;-1)$.\\
		Vì $A B C D$ là hình  thang nên $\overrightarrow{A B}$ và $\overrightarrow{D C}$ cùng hướng, nghĩa là
		 \allowdisplaybreaks
		\begin{eqnarray*}
			\exists k>0\colon \overrightarrow{D C}=k \overrightarrow{A B} 
			&\Leftrightarrow&\heva{&6-a=k \\& 1-b=-2 k \\&  -c=-2 k}\\ &\Leftrightarrow&\heva{&a=6-k \\& b=1+2 k \\& c=2 k.} \Rightarrow D(6-k;1+2 k;2 k).\\
			\left[\overrightarrow{A B}, \overrightarrow{A C}\right]=(0;-9;9)  &\Rightarrow& S_{\triangle A B C}=\dfrac{1}{2} \cdot\left|\left[\overrightarrow{A B}, \overrightarrow{A C}\right]\right|=\dfrac{1}{2} \cdot \sqrt{(-9)^2+9^2}=\dfrac{9 \sqrt{2}}{2}.\\
			\overrightarrow{A D}&=&(5-k;2 k-1;2 k-1),\\ 
			\left[\overrightarrow{A D}, \overrightarrow{A C}\right]&=&(0;9 k;-9 k) \\\Rightarrow S_{\triangle A D C}&=&\dfrac{1}{2} \cdot \left|\left[\overrightarrow{A D}, \overrightarrow{A C}\right]\right|=\dfrac{1}{2} \cdot \sqrt{(9 k)^2+(-9 k)^2}=\dfrac{9 k \sqrt{2}}{2}.\\ 
			S_{A B C D}&=&S_{\triangle A B C}+S_{\triangle A D C}=\dfrac{9 \sqrt{2}}{2}(1+k)=6 \sqrt{2}  \\\Leftrightarrow 1+k&=&\dfrac{4}{3} \Rightarrow k=\dfrac{1}{3}.
		\end{eqnarray*}
		Vậy $D\left(\dfrac{17}{3};\dfrac{5}{3};\dfrac{2}{3}\right)  \Rightarrow a+b+c=\dfrac{17+5+2}{3}=8$.
	 }
\end{ex}


%%%%%%%%%%%%%%%% ĐÚNG/SAI %%%%%%%%%%%%%%%%

% nguồn: Lop_12_Hinh_hoc_Chuong_2
\begin{ex}%[2H2H2-5]%[ma:872ecbaa]
	Trong hệ trục $Oxyz$, cho $3$ điểm $A\left(1;0;0\right)$,  $B\left(0;0;1\right)$, $C\left(2;1;1\right)$.
	 \choiceTF
		{\True Diện tích của  tam giác $ABC$ bằng $\dfrac{\sqrt{6}}{2}$ (đvdt)}
		 {Gọi $D\left(x;y;z\right)$ sao cho tứ giác $ABCD$ là một hình bình hành khi đó $x+y+z=3$}
		{\True Độ dài  đường cao của tam giác $ABC$ hạ từ $A$ bằng $AH=\dfrac{\sqrt{30}}{5}$ (đơn vị độ dài)}
		{Thể  tích của khối chóp $SABCD$ với đỉnh $S\left(0;3;4\right)$ bằng $2$ (đvtt)}
	 \loigiai{
		\begin{itemchoice}
			 \itemch Ta có $\overrightarrow{AB}=\left(-1;0;1\right),\overrightarrow{AC}=\left(1;1;1\right)$.\\
			Tính $[\overrightarrow{AB},\overrightarrow{AC}]=\left(\left|\begin{array}{ll} 0& 1\\ 1& 1
		\end{array}\right|;\left|\begin{array}{ll} 1& 1\\ {-1} & 1
	\end{array}\right|;\left|\begin{array}{ll}{-1} & 0\\ 1& 1
	\end{array}\right|\right)=\left(-1;2;-1\right)\ne \vec0$.\\
	Do đó $2$ vectơ $\overrightarrow{AB}$ và $\overrightarrow{AC}$ không cùng phương. Vậy $A$, $B$, $C$ là $3$ đỉnh của  một tam giác.\\
	Diện tích tam giác  $ABC$ $S_{\triangle ABC}=\dfrac{1}{2} \left|[\overrightarrow{AB},\overrightarrow{AC}]\right|=\dfrac{1}{2} \sqrt{\left(-1\right)^2+2^2+\left(-1\right)^2}=\dfrac{\sqrt{6}}{2}$ (đvdt).
	\itemch $ABCD$ là hình bình hành khi và chỉ khi $\overrightarrow{AD}=\overrightarrow{BC}$.\\ 
	Gọi $D\left(x;y;z\right)$ ta  có $\overrightarrow{AD}=\left(x-1;y;z\right);\overrightarrow{BC}=\left(2;1;0\right)$.\\
	Vậy $\overrightarrow{AD}=\overrightarrow{BC}\Leftrightarrow \heva{&x-1&=2 \\ &y&=1 \\ &z&=0}\Leftrightarrow \heva{&x=3 \\ &y=1 \\  &z=0}$. Vậy $x+y+z=4$.
	\itemch Diện tích $\triangle ABC=\dfrac{1}{2}  AH\cdot BC=\dfrac{\sqrt{6}}{2} \Leftrightarrow AH=\dfrac{\sqrt{6}}{BC}$.\\
	Ta có $BC=\sqrt{5}$ $\Leftrightarrow AH=\dfrac{\sqrt{30}}{5}$ (đơn vị dài). 
	\itemch Thể tích của khối chóp  $SABCD=V$.
	 Ta có $V=2V_{SABC}=\dfrac{1}{3} \left|[\overrightarrow{AB},\overrightarrow{AC}]\cdot \overrightarrow{AS}\right|$.\\
	Tính  $\overrightarrow{AS}=(-1;3;4)$ do kết quả câu a) nên $\left[\overrightarrow{AB},\overrightarrow{AC}\right]\cdot\overrightarrow{AS}=1+6-4=3> 0$.\\
	Vậy $V=1$  (đvtt).
	\end{itemchoice}
} 
\end{ex}

% nguồn: Lop_12_Hinh_hoc_Chuong_2
\begin{ex}%[2H2H2-5]%[ma:598dfba0]
	Trong không gian $Oxyz$, với ba vectơ đơn vị $\overrightarrow{i}, \overrightarrow{j}, \overrightarrow{k}$, cho $A(1; 1; 2), B(2; -1;  0), \overrightarrow{u}=(1; -2; -1)$.
	\choiceTF 
		{\True Tích có hướng của hai vectơ  $\overrightarrow{i}$ và $\overrightarrow{j}$ là vectơ $\overrightarrow{k}$}
		{$[\overrightarrow{u}, \overrightarrow{i}] =  (0; 1; 2)$}
		{$[\overrightarrow{AB}, \overrightarrow{u}] = (6; 1; 0)$} 
		{\True $[\overrightarrow{OA}, \overrightarrow{OB}] = (2; 4; -3)$} 
	 \loigiai{
		\begin{itemchoice}
			\itemch $\heva{ 
			&\overrightarrow{i} = (1; 0;  0) \\
			&\overrightarrow{j} = (0; 1;  0)
			} \Rightarrow [\overrightarrow{i}, \overrightarrow{j}] = (0; 0; 1)  = \overrightarrow{k}.$
			\itemch $\heva{ 
			&\overrightarrow{u} = (1; -2;  -1) \\
			&\overrightarrow{i}  = (1; 0; 0)
			} \Rightarrow [\overrightarrow{u}, \overrightarrow{i}] = (0; -1; 2) \neq (0; 1; 2)$. 
			\itemch  $\heva{
			&\overrightarrow{AB} = (1; -2; -2)  \\
			&\overrightarrow{u} = (1;  -2; -1)
			}  \Rightarrow [\overrightarrow{AB}, \overrightarrow{u}] = (-2; -1; 0) \neq (6; 1; 0)$.
			\itemch $\heva{ 
			&\overrightarrow{OA} = (1; 1;  2) \\
			&\overrightarrow{OB} = (2;  -1; 0)
			} \Rightarrow  [\overrightarrow{OA}, \overrightarrow{OB}] = (2; 4; -3)$.
		\end{itemchoice}
	} 
\end{ex}

% nguồn: Lop_12_Hinh_hoc_Chuong_2
\begin{ex}%[2H2H2-5]%[ma:ddbaea4c]
	Trong không gian $Oxyz$, cho các điểm $A(1;-3;  2)$, $B(2; 4;-2)$ và $C'(3; 2;-2)$.
	 \choiceTF
		 {\True Trung điểm của đoạn thẳng $O B$ là $C(1; 2;-1)$}
		{\True Biết rằng tứ giác $A C C' A'$  là hình bình hành. Cao độ của điểm $A'$ là $z=1$}
		{Biết rằng điểm $B'$ là  đỉnh còn lại của hình lăng trụ $A B C . A' B' C'$. Khi đó tung độ của điểm $B'$ là $y=-3$}
		{\True Thể tích của  khối lăng trụ $A B C . A' B' C'$ bằng $\dfrac{7}{2}$}
	\loigiai{ 
		\begin{itemchoice}
			 \itemch \textbf{Đúng}.
			Toạ độ  trung điểm $C\left(x_C; y_C; z_C\right)$ của đoạn thẳng $O B$ là
			$$\heva{&x_C=\dfrac{x_B+x_O}{2}  \\ &y_C=\dfrac{y_B+y_O}{2} \\& z_C=\dfrac{z_B+z_O}{2}} \Rightarrow\heva{&x_C=1 \\ &y_C=2 \\ &z_C=-1 .}$$
			Suy ra $C(1; 2;-1)$. 
			 \itemch \textbf{Đúng}.
			Vì $A$, $C$, $C'$ không thẳng hàng nên để tứ giác $A C C' A'$ là hình bình hành thì $$\overrightarrow{A A'}=\overrightarrow{C  C'}.$$
			Suy ra $A'(3;-3; 1)$, do đó cao  độ của điểm $A'$ là $z=1$.
			\itemch  \textbf{Sai}.
			Vì điểm $B'$ là đỉnh còn lại của hình lăng trụ $A B C.A'  B' C'$ nên $\overrightarrow{B B'}=\overrightarrow{C C'}$.\\
			Suy ra $B'(4; 4;-3)$. Do  đó tung độ của điểm $B'$ là $y=4$.
			\itemch \textbf{Đúng}. 
			Ta có $V_{A B C. A' B' C'}=S_{A B C} \cdot  \mathrm{d}\left(A',(A B C)\right)$.
			 $$
			\left[\overrightarrow{A B}, \overrightarrow{A C}\right]=(-1; 3; 5) \text {nên  } S_{A B C}=\dfrac{1}{2}\left|\left[\overrightarrow{A B}, \overrightarrow{A C}\right]\right|=\dfrac{\sqrt{35}}{2}.
			 $$
			Phương trình mặt phẳng đi qua ba  điểm $A$, $B$, $C$ là $x-3 y-5 z=0$. Khi đó
			$$ 
			\mathrm{d}\left(A',(A B C)\right)=\dfrac{|3-3 \cdot(-3)-5 \cdot 1|}{\sqrt{35}}=\dfrac{7  \sqrt{35}}{35}.$$
			Vậy $V_{A  B C . A' B' C'}=\dfrac{7}{2}$.
		\end{itemchoice}
	} 
\end{ex}

% nguồn: Lop_12_Hinh_hoc_Chuong_2
\begin{ex}%[2H2V2-5]%[ma:a00c3992]
	Trong không gian $Oxyz$, cho ba điểm  $A(3;-2;5)$, $B(2;1;-3)$ và $C(-1;2; 1)$.
	 \choiceTF
		{\True Hình chiếu vuông góc của điểm $A$ trên trục $Oz$ là $A'(0;0;5)$} 
		{\True Hình chiếu vuông góc của điểm $B$ trên mặt phẳng $(Oyz)$ là  $B'(2;0;-3)$}
		{Biết rằng $E$ là điểm nằm trên trục $Ox$ và $F$  là nằm điểm trên mặt phẳng $(Oyz)$ sao cho $C$ là trung điểm của đoạn thẳng $EF$. Khi đó $EF>5$}
		{\True Gọi $I(a;b;c)$ là  trực tâm của tam giác $ABC$. Khi đó $a+b+c=2$}
	\loigiai{ 
		\begin{itemchoice}
			\itemch  Đúng. Hình chiếu vuông góc của điểm $A\left(x_{A};y_{A};z_{A}\right)$ lên trục $Oz$ là $A'\left(0;0;z_{A}\right)$.
			\itemch Đúng. Hình chiếu vuông góc của điểm $B\left(x_{B};y_{B};z_{B}\right)$  trên mặt phẳng $(Oyz)$ là $B'\left(0;y_{B};z_{B}\right)$.
			\itemch Sai. Gọi $E\left(x_{E};0;0\right) \in Ox$ và $F\left(0;y_{F};z_{F}\right) \in(Oyz)$  sao cho $C$ là trung điểm của đoạn thẳng $EF$.
			Khi  đó
			$$ 
			\heva{ 
			& {\dfrac {x _ {E} + x _  {F}}{2} = x _ {C}} \\
			& {\dfrac {y _ {E} + y _ {F}}{2}  = y _ {C}} \\
			& {\dfrac {z _ {E} + z  _ {F}}{2} = z _ {C}} }
			\Rightarrow  \heva{
			& {\dfrac {x _  {E} + 0}{2} = - 1} \\
			& {\dfrac {0 + y _  {F}}{2} = 2} \\
			& {\dfrac {0 + z _  {F}}{2} = 1}}
			\Rightarrow  \heva{
			 & {x _ {E} = - 2} \\
			& {y _  {F} = 4} \\
			& {z _ {F} =  2}} \Rightarrow \heva{
			&E(-2;0;0) \\ 
			 & F(0;4;2) }
			$$ 
			Suy ra $EF=2 \sqrt{6}<5$. 
			\itemch  Đúng.\\
			$I(a;b;c)$ là trực tâm của tam giác $ABC$  nên $\heva{
			& A I \perp  BC \\
			& B  I \perp AC \\
			&  I \in(ABC) }\Rightarrow\heva{
			& \overrightarrow{AI} \cdot \overrightarrow{BC}=0  \\
			& \overrightarrow{BI}  \cdot \overrightarrow{AC}=0 \\
			&  {\left[\overrightarrow{AB}, \overrightarrow{AC}\right] \cdot \overrightarrow{AI}=\overrightarrow{0}} } $
			Ta  có $\overrightarrow{AI}=(a-3;b+2;c-5), \overrightarrow{BI}=(a-2;b-1;c+3)$,\\
			$\overrightarrow{BC}=(-3;1;4)$, 
			 $\overrightarrow{AC}=(-4;4;-4), \overrightarrow{AB}=(-1;3;-8)$.\\
			Khi  đó $\heva{
			 & -3(a-3)+1(b+2)+4(c-5)=0 \\
			&  -4(a-2)+4(b-1)-4(c+3)=0 \\
			&  20(a-3)+28(b+2)+8(c-5)=0 }\Leftrightarrow\heva{
			&  3a-b-4c=-9 \\
			& a-b+c=-2 \\ 
			& 5a+7b+2c=11}  \Leftrightarrow \heva{
			 & a=-1 \\
			 & b=2 \\
			& c=1.  }$\\
			Suy  ra $I(-1;2;1)$. Vậy $a+b+c=-1+2+1=2$.
	\end{itemchoice}}
\end{ex}


%%%%%%%%%%%%%%%% TRẢ LỜI NGẮN %%%%%%%%%%%%%%%%

% nguồn: Lop_12_Hinh_hoc_Chuong_2
\begin{bt}%[2H2H2-5]%[ma:c12241cd]
	Trong không gian với hệ trục tọa độ $Oxyz$, cho tứ diện $ABCD$ có $A(2;-1;1)$, $B(3;0;-1)$, $C(2;-1;3)$, $D \in O y$ và thể tích tứ diện $ABCD$ bằng $5$. Tính tổng tung độ  của các điểm $D$ thỏa mãn yêu cầu bài toán.\shortans{-6}
	 \loigiai{
		 Do $D \in O y \Rightarrow D(0;y;0)$.
		Khi đó  $$\overrightarrow{D A}=(2;-1-y;1), \overrightarrow{D B}=(3;-y;-1), \overrightarrow{D C}=(2;-1-y;3).$$
		Ta có $$\left[\overrightarrow{D A},  \overrightarrow{D B}\right]=(1+2 y;5;y+3)\Rightarrow\left[\overrightarrow{D A}, \overrightarrow{D B}\right] \cdot \overrightarrow{D C}=2+4 y-5-5 y+3 y+9=2 y+6.$$
		Và $$V_{ABCD}=\dfrac{1}{6}\left|\left[\overrightarrow{D  A}, \overrightarrow{D B}\right] \cdot \overrightarrow{D C}\right|=5 \Leftrightarrow\hoac{&2 y+6=30 \\& 2 y+6=-30} \Leftrightarrow\hoac{&y=12 \\& y=-18.}$$
		Vậy $y_1+y_2=12-18=-6$. 
	} 
\end{bt}

% nguồn: Lop_12_Hinh_hoc_Chuong_2
\begin{bt}%[2H2H2-5]%[ma:55beb32e]
	Trong không gian với hệ tọa độ $Oxyz$, cho hai điểm $C(4;0;0)$ và $B(2;0;0)$. Điểm $M\left(0;y_M;0\right)$ thuộc trục tung sao cho diện tích tam giác  $MBC$ bằng $3$. Có bao nhiêu giá trị nguyên của $y_M$?
	 \shortans{2}
	\loigiai{Ta có $\overrightarrow{B M}=\left(-2;y_{M};0\right)$; $\overrightarrow{B  C}=(2;0;0)$.\\
		Khi đó $\left[\overrightarrow{B M},  \overrightarrow{B C}\right]=\left(0;0;-2 y_{M}\right)$.\\
		 Diện tích tam giác $MBC$ là $$S_{M B C}=\dfrac{1}{2}\left|\left[\overrightarrow{B M}, \overrightarrow{B C}\right]\right| \Leftrightarrow \dfrac{1}{2} \sqrt{4 y_{M}^2}=3 \Leftrightarrow\left|y_M\right|=3 \Leftrightarrow y_M= \pm 3.$$
		Vậy có $2$ giá trị nguyên của $y_M$. 
	} 
\end{bt}

% nguồn: Lop_12_Hinh_hoc_Chuong_2
\begin{bt}%[2H2C2-5]%[ma:af405054]
	Trong không gian với hệ trục tọa độ $Oxyz$, một tấm bằng đồng chất có dạng hình vuông $ABCD$ tâm $O'$ được treo nghiêng bởi 4 sợi dây $IA$, $IB$, $IC$, $ID$ gắn cố định tại điểm $I(0;0;9)$ như hình vẽ. Biết điểm $O'(0;0;5)$, $A(-\sqrt{2}; \sqrt{2}; 5)$, điểm $D(a;b;c)$ với $a>0$, $c>5$ và góc $\widehat{O'ID}$ đạt giá trị nhỏ nhất. Gọi  $\vec{F}_1$, $\vec{F}_2$, $\vec{F}_3$, $\vec{F}_4$ lần lượt là lực căng của dây $IB$, $IC$, $ID$, $IA$, biết $|\vec{F}_2| = 50N$ và $|\vec{F}_4| = 40N$.
	\begin{center}
		\begin{tikzpicture}[line join = round,line cap = round,>=stealth,font=\footnotesize,scale=1]
				\coordinate (A)  at (0,0);
				\coordinate  (B) at ($(A)+(5,0)$);
				\coordinate (D) at  ($(A)+(-140:2.7)$);
				\coordinate (C) at  ($(B)+(D)-(A)$);
				\path[name path=AC] (A) -- (C); 
				\path[name path=BD] (B)  -- (D);
				\path[name intersections={of=AC  and BD,by=O'}];
				 \coordinate (I) at ($(O')+(0,4)$);
				\draw  (I) -- (B) -- (C) -- (D) -- cycle;
				\fill[green] (A) -- (B) -- (C)  -- (D) -- cycle;
				\draw (I) -- (C); 
				\draw[dashed] (A)  -- (I)
					(A) --  (B)
					(A)  -- (D)
					 (A) -- (C)
					(B) --  (D);
				 \foreach \pt/\pos in {A/left,B/right,C/right,D/left,O'/below,I/above}
					{ 
				\fill (\pt)  circle (2pt);
				\node[\pos=2pt]  at (\pt) {$\pt$};
				} 
			\end{tikzpicture}
	\end{center}
	Tính khối lượng $m$(kg) của tấm bằng (làm tròn kết quả đến hàng đơn vị) (cho gia tốc trọng trường $g=9{,}8$ m/s$^2$)  và trọng lượng của tấm bằng được tính bằng công thức $P=mg$).
	\shortans{15} 
	 \loigiai{
		\begin{center}
			\begin{tikzpicture}[line join = round,line cap = round,>=stealth,font=\footnotesize,scale=1]
					 \coordinate (A) at (0,0);
					\coordinate (B)  at ($(A)+(5,0)$);
					\coordinate (D) at ($(A)+(-140:2.7)$); 
					\coordinate (C)  at ($(B)+(D)-(A)$);
					\path[name path=AC]  (A) -- (C);
					\path[name path=BD] (B) --  (D);
					\path[name intersections={of=AC  and BD,by=O'}];
					\coordinate (I) at  ($(O')+(0,4)$);
					\draw (I)  -- (B) -- (C) -- (D) -- cycle;
					\draw (I)  -- (C);
					\draw[dashed]  (A) -- (I)
						 (A) -- (B)
						(A) -- (D) 
						(A) --  (C)
						(B)  -- (D)
						(I)--(O'); 
					\foreach \pt/\pos  in {A/left,B/right,C/right,D/left,O'/below,I/above}
						{ 
					\fill (\pt)  circle (2pt);
					 \node[\pos=2pt] at (\pt) {$\pt$};
					 }
				\end{tikzpicture}
		\end{center}
		Ta có  $O'A=2$, $AC=4$, $AB=2\sqrt{2}$, $O'D=2$, $O'I=4$, $IA=2\sqrt{5}$.\\
		Vì $O'$ là trung  điểm $AC \Rightarrow C(\sqrt{2};-\sqrt{2};5) \Rightarrow IC=2\sqrt{5}$.\\
		Vì điểm $O'$ là tâm của hình vuông $ABCD \Rightarrow \vec{O'D}  \perp \vec{O'A} \Rightarrow \vec{O'D}\cdot \vec{O'A}=0 \Rightarrow a=b$.\\
		Góc $\widehat{O'ID}$ đạt giá trị lớn nhất $\Leftrightarrow \cos \widehat{O'ID}$ nhỏ  nhất.\\
		Áp dụng định lý cô-sintrong tam giác $O'ID$  ta có
		$$\cos \widehat{O'ID} = \dfrac{ID^2+O'I^2-O'D^2}{2ID.O'I} = \dfrac{12+ID^2}{8ID} = \dfrac{ID}{8} + \dfrac{3}{2ID} \ge  2\sqrt{\dfrac{3}{16}} = \dfrac{\sqrt{3}}{2}.$$
		Suy  ra $\widehat{O'ID} = 30^\circ$ và $ID=2\sqrt{3}$.\\
		Vì $\heva{&ID=2\sqrt{3} \\ &O'D=2} \Leftrightarrow \heva{&a^2+b^2+(c-9)^2=12 \\ &a^2+b^2+(c-5)^2=4} \Leftrightarrow \heva{&(c-9)^2-(c-5)^2=8  \\&a^2+b^2+(c-5)^2=4}\Leftrightarrow\heva{&c=6 \\&a=b=\dfrac{\sqrt{6}}{2}}$ (do $a>0$, $c>5$).\\
		Suy ra $D\left(\dfrac{\sqrt{6}}{2};\dfrac{\sqrt{6}}{2};6\right) \Rightarrow  B\left(-\dfrac{\sqrt{6}}{2};-\dfrac{\sqrt{6}}{2};4\right) \Rightarrow ID=2\sqrt{3}$, $IB=2\sqrt{7}$.\\
		Vì $IA=IC=2\sqrt{5}$ nên suy ra  tam giác $IAC$ cân tại $I$ nên $\left| \vec{F}_2\right| =\left| \vec{F}_4\right| = 40N$.
		\begin{center}
			\begin{tikzpicture}[line join = round,line cap = round,scale=1]
					\path 
						 (0,0) coordinate (A)
						(6,0) coordinate (C) 
						(4,3) coordinate (I) 
						($(A)!0.5!(C)$) coordinate  (O')
						($(I)!0.25!(A)$) coordinate (A_1) 
						 ($(I)!0.25!(C)$) coordinate (C_1)
						($(A_1)+(C_1)-(I)$) coordinate  (i)
						($(A_1)!0.5!(C_1)$)  coordinate (H);
					 \draw (I)--(A)--(C)--cycle;
					\draw  (A_1)--(i)--(C_1);
					\draw (I)--(O') 
						 (A_1)--(C_1);
					\foreach \x/\y in {I/above,A/left,C/right,O'/below,A_1/left,C_1/right,H/above right}{ 
						 \node at (\x)[\y]{$\x$};
						 \fill[red] (\x) circle (1.5pt);}
				\end{tikzpicture}
		\end{center}
		Dựng $\vec{IA_1}=\vec{F}_2+\vec{F}_4$ và $H$ là trung điểm $A_1C_1$.\\ 
		 Tam giác $IAO'$ vuông tại $O' \Rightarrow \cos \widehat{AIO'} = \dfrac{IO'}{IA} = \dfrac{4}{2\sqrt{5}} = \dfrac{2}{\sqrt{5}}$.\\
		$\Rightarrow \left|\vec{F_2}+\vec{F_4}\right| = \left|\vec{IA}+\vec{IC}\right| = 2|\vec{IH}|=2|\vec{IA}|\cos \widehat{AIO'} = 2\cdot \left(2\sqrt{5}\right)\cdot  \dfrac{2}{\sqrt{5}} = 8$.\\
		\begin{center}
			\begin{tikzpicture}[line join = round,line cap = round,scale=1]
					\path 
						(0,0)  coordinate (B)
						(6,0) coordinate (D) 
						 (4,3) coordinate (I)
						 ($(B)!0.5!(D)$) coordinate (O')
						 ($(I)!0.25!(B)$) coordinate (B_1)
						($(I)!0.25!(D)$)  coordinate (D_1)
						 ($(B_1)+(D_1)-(I)$) coordinate (K)
						($(B_1)!0.5!(D_1)$) coordinate  (H);
					\draw (I)--(B)--(D)--cycle; 
					\draw  (I)--(O')
						(B_1)--(K)--(D_1); 
					\draw[->] (I)--(B_1); 
					 \draw[->](I)--(D_1);
					 \draw[->](I)--(K);
					 \foreach \x/\y in {I/above,B/left,D/right,O'/below,B_1/left,D_1/right,K/right}{
						\node at  (\x)[\y]{$\x$};
						\fill[red] (\x) circle (1.5pt);} 
				\end{tikzpicture}
		\end{center}
		Dựng $\vec{IB_1}=\vec{F_1}+\vec{F}_3$ và $K$ là điểm thứ 4  của hình bình hành $IB_1KD_1$.
		 $\Rightarrow \left|\vec{F}_1+\vec{F}_3\right| = |\vec{IB}|+|\vec{ID}| = |\vec{IK}|$.\\
		Áp  dụng định lí côsintrong tam giác $IBD$ ta có:
		$$\cos \widehat{B_1IO'} = \dfrac{IB^2+O'I^2-O'B^2}{2IB.O'I} = \dfrac{101+16-4}{2\sqrt{101}.4} = \dfrac{113}{8\sqrt{101}} \approx 1{,}40625  \Rightarrow \widehat{B_1IK} \approx 19^\circ.$$
		Mặt khác $|\vec{IB_1}|=|\vec{DIO'}|=30^\circ  \Rightarrow \widehat{IB_1K} \approx 131^\circ$.\\
		Áp  dụng định lí sintrong tam giác $IB_1K$ ta có:
		$$\dfrac{IK}{\sin 131^\circ}  = \dfrac{IB_1}{\sin 30^\circ} \Rightarrow IK = \dfrac{IB_1 \sin 131^\circ}{\sin 30^\circ} = \dfrac{5 \cdot 50 \cdot \sin 131^\circ}{0{,}5} \approx 75{,}47.$$
		Khi vật cân bằng 
		\begin{eqnarray*}
			&&\vec{F}_1+\vec{F}_2+\vec{F}_3+\vec{F}_4+\vec{P} = \vec{0}\\ 
			&\Leftrightarrow& \vec{F}_1+\vec{F}_2+\vec{F}_3+\vec{F}_4 =  - \vec{P}\\
			&\Leftrightarrow& \left|2\vec{IH}+\vec{IK}\right|  =mg\\
			&\Leftrightarrow& 2\left|\vec{IH}\right|+\left|\vec{IK}\right| =  mg\\
			&\Leftrightarrow& 32\sqrt{5}+75{,}47  = mg\\
			&\Leftrightarrow& m \approx 15 \text{  ~kg}.
		\end{eqnarray*}
	} 
\end{bt}


%%%%%%%%%%%%%%%% TỰ LUẬN %%%%%%%%%%%%%%%%

% nguồn: Lop_12_Hinh_hoc_Chuong_2
\begin{bt}%[2H2H2-5]%[ma:fdee3069]
	Trong không gian $Oxyz$, cho hai vectơ $\overrightarrow{m}=(4;3;1)$, $\overrightarrow{n}=(0;0;1)$. Gọi $\overrightarrow{p}$ là vectơ cùng  hướng với $\left[\overrightarrow{m}, \overrightarrow{n}\right]$ (tích có hướng của hai vectơ $\overrightarrow{m}$ và $\overrightarrow{n}$). Biết $\left|\overrightarrow{p}\right| =15$, tìm tọa độ vectơ $\overrightarrow{p}$.
	 \loigiai{
		Ta  có $\left[\overrightarrow{m},\overrightarrow{n}\right]=(3;-4;0)$.\\
		Do $\overrightarrow{p}$ là vectơ cùng hướng với $\left[\overrightarrow{m},\overrightarrow{n}\right] $ nên $\overrightarrow{p}=k\left[\overrightarrow{m},\overrightarrow{n}\right]$, $k>0$.\\ 
		Mặt khác, $\left|\overrightarrow{p}\right|=15 \Leftrightarrow k \cdot\left|\left[\overrightarrow{m},\overrightarrow{n}\right]\right|=15 \Leftrightarrow k \cdot 5=15  \Leftrightarrow k=3$.\\
		 Vậy $\overrightarrow{p}=(9;-12;0)$.
	} 
\end{bt}
